\section{质量与多样性是一笔交易}
\label{sec:14-Deep-Learning/07-Generative-Models-Unified/04}

评审会上摆着两个模型的数：A 的 FID 是 \(12.4\)，B 是 \(18.7\)。按指标选 A，上线。两周后运营那边转来一堆反馈，说 A 生成的图``翻来覆去就那么几种''，反而是 B 虽然偶尔出废图，用户愿意多点几次重新生成。有人提议再调调把 A 的 FID 压到 \(10\) 以下。

这个提议解决不了问题，因为分歧根本不在那个数的大小上。会上真正在争的是两件事：生成出来的东西像不像真的，以及真实世界里各种各样的东西模型能不能生成出来。这两件可以互相牺牲——把模型锁死在最有把握的那几种输出上，第一件立刻变好，第二件立刻变差。一个数不可能同时报告两个可以对着换的量，它只能按自己内定的比例把两者混成一个分数，而那个比例不是你的业务定的。

\subsection{把一个数拆成两个}

拆法可以直接从分布的支撑上读出来。生成分布 \(p_\theta\)，数据分布 \(p_{\text{data}}\)，定义

\[
\begin{aligned}
\text{precision}&=P\!\left(x\in\operatorname{supp}(p_{\text{data}})\ \middle|\ x\sim p_\theta\right),\\
\text{recall}&=P\!\left(x\in\operatorname{supp}(p_\theta)\ \middle|\ x\sim p_{\text{data}}\right).
\end{aligned}
\]

第一行问：我生成的东西里，有多大比例是真实世界里可能出现的？这是质量。第二行问：真实世界里的东西，有多大比例是我够得到的？这是多样性。两个概率的条件方向相反，量的是两件不同的事。实际估计时支撑不可能精确知道，做法是把样本映到某个特征空间，用每个点到它第 \(k\) 个近邻的距离撑起一个小球，并集当作支撑的近似——这层近似的可靠性后面还要回来算账。

现在把第 01 篇那个不对称接上来。前向 KL 在 \(p_{\text{data}}>0\) 而 \(p_\theta\approx0\) 处受重罚，那正是 recall 掉下去的地方；它不罚 \(p_\theta\) 跑到数据之外，而那正是 precision 掉下去的地方。反向 KL 的账恰好反过来。两个 KL 方向不是``一个准一点一个糙一点''，是各自把这两个指标中的一个推到极端而不管另一个。JS 因为对称，两侧都罚一点，落在中间。

\begin{center}
\begin{tikzpicture}[x=4.6cm,y=3.9cm,>=stealth]
  \draw[->] (-0.03,0) -- (1.12,0) node[right,font=\small] {recall（多样性）};
  \draw[->] (0,-0.03) -- (0,1.12) node[above,font=\small] {precision（质量）};
  % 弱模型的权衡曲线
  \draw[blue!70!black,thick] plot[smooth] coordinates
    {(0.08,0.93) (0.25,0.88) (0.45,0.78) (0.62,0.63) (0.75,0.45) (0.83,0.25) (0.88,0.08)};
  % 更强模型：整条曲线外移
  \draw[green!45!black,thick,dashed] plot[smooth] coordinates
    {(0.12,0.99) (0.35,0.96) (0.58,0.90) (0.75,0.80) (0.86,0.62) (0.93,0.38) (0.96,0.14)};
  % 三个工作点
  \fill[blue!70!black] (0.25,0.88) circle (2.3pt);
  \fill[blue!70!black] (0.62,0.63) circle (2.3pt);
  \fill[blue!70!black] (0.83,0.25) circle (2.3pt);
  \node[font=\small,anchor=north east] at (0.235,0.855) {低温};
  \node[font=\small,above right] at (0.645,0.645) {默认};
  \node[font=\small,anchor=north east] at (0.815,0.215) {高温};
  \node[blue!70!black,font=\small] at (0.26,0.52) {同一组权重};
  \node[green!35!black,font=\small] at (0.72,0.99) {换一个更强的模型};
  \draw[->,gray!75] (0.50,0.745) -- (0.62,0.865);
\end{tikzpicture}

\emph{调温度是在同一条曲线上滑动，模型变强是整条曲线往外挪。单个 FID 数值分不清这两件事，而它们对应完全不同的行动。}
\end{center}

\subsection{似然和感知质量可以往两个方向背离}

第一件容易踩的事：把对数似然当成质量的代理。它不是，而且不是``大体一致偶有例外''，是两个方向都能坏。

\begin{example}
设 \(q\) 是纯噪声分布，构造混合模型

\[
p_\theta=0.99\,q+0.01\,p_{\text{data}}.
\]

从它采样，\(99\%\) 的输出是噪声。可对任意 \(x\)，

\[
\log p_\theta(x)\ \ge\ \log\left(0.01\,p_{\text{data}}(x)\right)=\log p_{\text{data}}(x)-\log 100 .
\]

取期望，这个模型的平均对数似然最多比真分布低 \(\log 100\approx4.6\) 奈特。报告里通常按维度折算：\(d=3072\) 的图像上是 \(4.6/3072\approx0.0015\) 奈特每维，落在小数点后第三位，任何一次超参数调整带来的波动都比它大。指标上看不出任何异常，采样出来九成九是雪花。
\end{example}

反方向同样成立。一个只学会了单一模式但把它学得极好的模型，生成的每张图都无可挑剔，可它对其余模式给出零密度，对数似然是 \(-\infty\)。上一篇那条只贴住一个峰的曲线就是它。

根子在于两个量的积分方向不同。对数似然是在 \(p_{\text{data}}\) 下取平均——它问的是``真实样本在模型眼里有多可信''；人看到的样本来自 \(p_\theta\)——问的是``模型吐出来的东西有多真''。高维空间里 \(p_\theta\) 可以把绝大部分质量放在没人会去看的区域，同时在数据附近留下刚好够用的一点点密度。似然只看后面那点，眼睛只看前面那堆。上一篇提到的那个经验观察——在 CIFAR-10 上训练的流给 SVHN 图片更高的似然——是同一件事在真实模型上的表现。

\subsection{温度和截断是在曲线上滑动，不是在改进模型}

第二件容易踩的事：把采样设置的调整读成模型能力的提升。

自回归模型里的温度是对每步的条件分布做 \(p_\tau(x_t\given x_{<t})\propto p(x_t\given x_{<t})^{1/\tau}\)。\(\tau<1\) 把分布削尖，质量从长尾挪向高概率的少数选项，输出更``稳''也更单调；\(\tau>1\) 反过来。截断更直接：top-\(k\) 与 nucleus 采样把低概率的尾巴整段切掉再重新归一化；GAN 那边的 truncation trick 把噪声 \(z\) 限制在一个半径内，等于砍掉潜空间的边缘区域。

这些操作有一个共同点：一个参数都没有变。同一组权重，改一个采样常数，precision 和 recall 就同向反向地各动一格——正是图里蓝色曲线上那三个点。所以``降温之后 FID 从 \(18\) 掉到 \(12\)''这句话不能读成模型变好了，它说明的是 FID 在这条权衡曲线上偏好靠近高 precision 的那一段。要比较两个模型，得比较两条曲线，而不是各自随手挑一个工作点上的数。可复现的做法是扫一遍温度，把整条曲线画出来，再看有没有一条整体压住另一条——图里绿色虚线那种情形才叫模型更强。

\subsection{FID 依赖三样东西，每一样都会影响结论}

第三件容易踩的事：把 FID 当成一个客观的距离。

它的计算是：把两组样本送进一个预训练网络（通常是 ImageNet 上训练的 Inception），取某一层的激活，分别估出均值与协方差，再算两个高斯之间的 Fréchet 距离

\[
\text{FID}=\norm{\mu_1-\mu_2}^2+\trace\!\left(\Sigma_1+\Sigma_2-2\left(\Sigma_1\Sigma_2\right)^{1/2}\right).
\]

三处依赖都写在这个流程里。

激活分布被当成高斯处理，可它不是。这一步把两个分布的差异压缩成前两阶矩的差异，高阶矩上的任何区别——多峰、偏斜、重尾——一律看不见。两个明显不同的分布只要均值协方差对上，FID 就是零。

特征提取器决定了``像''的含义。FID 量的是两个分布在\emph{那个}特征空间里有多接近。Inception 的表征是为 ImageNet 的一千类物体训练出来的，对自然图像里的物体身份、纹理、构图很敏感。换成医学影像、卫星图、字体渲染这类域外数据，它未必编码了业务真正在意的差异，换一个特征提取器排名可能就变。所以两篇论文的 FID 只有在特征提取器与实现细节完全一致时才可比。

样本数带来的是系统性偏差，不是噪声。\(\Sigma\) 是 \(2048\times2048\) 的协方差矩阵，用 \(n\) 个样本估计；\(n\) 不够大时协方差被系统性地估偏，迹项因此带一个正的偏移，FID 随 \(n\) 增大而单调下降，通常要到几万个样本才趋于稳定。这是有限样本偏差，多跑几次取平均消不掉。用 \(10^4\) 个样本算出的 FID 和用 \(5\times10^4\) 个算出的 FID 不能放在一张表里比。

能做的事很朴素：报告里至少给两维（precision 与 recall，或它们的稳健变体），把样本数、特征提取器版本、采样设置一并写上，比较模型时给出扫过温度的整条曲线。这些不能让评价变得客观，但至少让两次评价说的是同一件事。

\subsection{回到那个共同的困难}

四条路线现在都放回了同一个框架：都在让 \(p_\theta\) 靠近 \(p_{\text{data}}\)，区别只在用哪个散度量差距、以及怎样绕开算不出来的那一项。VAE 用下界绕开边缘似然，流用可逆性让归一化常数恒为一，GAN 干脆不写密度、把散度的估计外包给一个网络，三者各自赔进去的是模糊、架构受限、以及理论保证的缺失。

放到一起看，剩下的困难其实是同一个：训练开始时 \(p_\theta\) 和 \(p_{\text{data}}\) 离得很远，而所有散度在``离得远''的时候都不好用——JS 直接饱和成常数，两个 KL 方向则被迫在覆盖与贴合之间做极端取舍，于是有了模糊和模式坍缩。换一个更聪明的散度只能缓解，改不掉这个结构。

真正的出路是不要一次性匹配。给数据逐步加噪声，一步一步加到它变成纯高斯；这条路上相邻两个噪声尺度之间的分布差别很小，小到匹配它们只需要一个平方误差回归，既不需要归一化常数，也不需要判别器，更不会遇到支撑不重叠。把这些小步骤反着连起来，就得到一条从高斯走回数据的路径，而每一步都容易。下一个单元处理的正是这条路线：它仍然站在前向 KL 的方向上，但把一次跨越整个分布的困难匹配，换成了沿噪声尺度排开的很多次容易匹配。
